\section{One Kingship, Two Questions}
The first reading asks how a person becomes a subject who hears. Augustine's answer begins with grace: the kingdom's origin is not the world, and inward hearing follows the gift of belonging to truth. Jerome supplies the enduring human inheritance opened by the incarnate Son. The principal movement is from rescue to obedient faith and perseverance. Its distinctive temptation is self-sufficiency, whether expressed as confidence in one's own merit or in force that claims to manufacture the kingdom.

The second reading asks what universal lordship does for estranged creatures. Chrysostom's answer begins with the Creator who sustains all and gives himself to reconcile them; Bellarmine distinguishes rightful authority, willing conversion and the acknowledgment of judgment. The principal movement is from creative primacy through sacrifice into shared peace. Its distinctive temptation is to confess a universal kingdom while preserving private enmity or overlooking the helpless whom royal justice protects.

Their agreement does not erase their different starting points. Both locate the kingdom within the world without deriving it from worldly force. Both depend upon grace, sacrifice and a future not reducible to political success. Yet the first exposes a mistaken account of entry, while the second exposes a mistaken account of communion. The believer who listens humbly still has to learn what reconciliation demands; the believer who works for peace still receives the capacity to hear as a gift.

The genuine difference between Augustine and Bellarmine on Psalm 71 remains useful within that agreement. Augustine introduces Christ as the true Solomon; Bellarmine allows the opening prayer its immediate royal setting before tracing the fulfillment beyond it. The former brings the peacemaker directly into view; the latter shows how historical justice and Christological fulfillment belong together. Neither licenses the conclusion that the psalm's universal homage is already the willing obedience of every ruler.

The feast ends by asking for lasting reign with the King whose life has been given for his people. Hearing and peace converge there. A kingdom whose truth never reaches obedience has been reduced to information; a peace that leaves every chosen feud intact has been reduced to a greeting. The Mass gives both a deeper reality: persons made fit for the Son's inheritance, and a people learning to live from the reconciliation his Cross accomplishes.
